El presente trabajo abordó la evaluación objetiva de cadenas de preprocesamiento para la generación de conjuntos de datos del habla a partir de material \textit{in-the-wild}, articulando dos paradigmas de evaluación complementarios: un conjunto de métricas acústicas objetivas, sintetizadas en una métrica compuesta de optimización, y un método basado en la estimación de densidad mediante modelos de difusión probabilística. La integración de ambas perspectivas permitió no solo caracterizar el comportamiento de las distintas configuraciones de la cadena, sino también exponer las limitaciones y los sesgos inherentes a cada criterio de evaluación.

El hallazgo de mayor relevancia se desprende de la divergencia sistemática entre ambas metodologías. Según el ordenamiento provisto por las métricas objetivas intrusivas, la configuración óptima recae de manera consistente en el uso de Demucs como algoritmo de realce, alcanzando el mínimo global de la métrica conjunta. Este comportamiento se explica porque métricas como PESQ o SI-SDR premian de forma casi lineal la ausencia de energía en las bandas donde se espera silencio, favoreciendo la agresividad con la que dicho algoritmo suprime el ruido de fondo. Sin embargo, desde la perspectiva de la estimación de densidad el ordenamiento se invierte drásticamente: Demucs obtiene los peores puntajes de verosimilitud, mientras que DeepFilterNet y la condición sin denoising dominan el ranking. Se concluye que esta contradicción no constituye un fallo metodológico, sino que evidencia que cada sistema conceptualiza la calidad de manera distinta. El procesamiento agresivo introduce huecos espectrales y discontinuidades topológicamente improbables en el habla natural, anomalías que el modelo de difusión detecta y penaliza por alejar la señal de la distribución de referencia, aun cuando las métricas objetivas las interpreten como una mejora.

En segundo lugar, la evaluación exploratoria sobre tareas de síntesis \textit{zero-shot} permitió contrastar la capacidad predictiva de ambos criterios frente a la calidad percibida del audio sintetizado. Se demostró que la métrica objetiva total fracasó al intentar predecir la calidad subjetiva estimada, llegando incluso a exhibir una tendencia inversa a la esperada y careciendo de significancia estadística. Este resultado se atribuye principalmente a la penalización asignada a la reducción de datos, factor irrelevante en escenarios de clonación de voz con referencias cortas. Por el contrario, la estimación de densidad reveló una correlación lineal positiva, fuerte y estadísticamente significativa con la calidad sintetizada para ambos sintetizadores evaluados. Se concluye, por tanto, que la verosimilitud computada mediante modelos de difusión emerge no solo como una medida teórica de distancia distribucional, sino como un indicador empíricamente válido para la selección de corpus destinados al entrenamiento de sistemas generativos.

Un tercer aporte se vincula con la sensibilidad arquitectónica observada entre los sintetizadores. Los resultados sugieren que la tolerancia de un modelo a los artefactos de preprocesamiento está condicionada por su arquitectura subyacente. El sintetizador autoregresivo basado en modelos de lenguaje, al operar sobre representaciones discretizadas del audio mediante codecs neuronales, exhibió una mayor resiliencia frente a ruidos residuales y distorsiones leves, ya que la cuantización en tokens acústicos actúa como un filtro de normalización intrínseco. En contraposición, el modelo basado en \textit{flow-matching} y difusión, cuyo objetivo matemático es el emparejamiento exacto de distribuciones, demostró ser hipersensible a la condición inicial, replicando fielmente las anomalías espectrales presentes en el audio de referencia. Se concluye que los modelos de difusión constituyen evaluadores más estrictos y transparentes de la calidad topológica real de un conjunto de datos, mientras que los modelos autoregresivos ofrecen mayor robustez en entornos de producción no controlados.

Finalmente, la comparación entre el corpus \textit{in-the-wild} y los conjuntos de datos profesionales validó la pertinencia del material recolectado de fuentes heterogéneas. Si bien el procesamiento de señal no logró igualar las condiciones de calidad alcanzables en grabaciones de estudio para todas las métricas, el corpus \textit{in-the-wild} exhibió una mayor variabilidad prosódica, cuantificada a través de la desviación estándar de la frecuencia fundamental, y una distribución de hablantes libre del sesgo hacia voces graves característico de los corpus profesionales. Se concluye que la incorporación de datos \textit{in-the-wild} resulta indispensable para capturar la diversidad real del fenómeno del habla y la representatividad dialectal del español de Argentina, atributos que favorecen la generalización de los modelos a condiciones reales de uso y que difícilmente pueden obtenerse mediante corpus controlados.